\section{Introduccion}
. presento el problema
. Empiezo a hablar como fue el aproach del problema
. planteo preguntas, que voy a responder despues

. hago un analisis del problema
. explico como llegamos a la mejor solucion y como pudimos mostrar que sea optimo (NO MOSTRAR, solo explicar, NO SPOILEAR ¿?)

\subsection{El problema}



El problema trataba de
Introducción y primeros años
Cuando Mateo nació, Sophia estaba muy contenta. Finalmente tendría un hermano con quien jugar. Sophi tenía 3 años cuando Mateo nació. Ya desde muy chicos, ella jugaba mucho con su hermano.

Pasaron los años, y fueron cambiando los juegos. Cuando Mateo cumplió 4 años, el padre de ambos le explicó un juego a Sophia: Se dispone una fila de 
n
n monedas, de diferentes valores. En cada turno, un jugador debe elegir alguna moneda. Pero no puede elegir cualquiera: sólo puede elegir o bien la primera de la fila, o bien la última. Al elegirla, la remueve de la fila, y le toca luego al otro jugador, quien debe elegir otra moneda siguiendo la misma regla. Siguen agarrando monedas hasta que no quede ninguna. Quien gane será quien tenga el mayor valor acumulado (por sumatoria).

El problema es que Mateo es aún pequeño para entender cómo funciona esto, por lo que Sophia debe elegir las monedas por él. Digamos, Mateo está “jugando”. Aquí surge otro problema: Sophia es muy competitiva. Será buena hermana, pero no se va a dejar ganar (consideremos que tiene 7 nada más). Todo lo contrario. En la primaria aprendió algunas cosas sobre algoritmos greedy, y lo va a aplicar.

Consigna
[Obligatorio] Hacer un análisis del problema, y proponer un algoritmo greedy que obtenga la solución óptima al problema planteado: Dados los 
n
n valores de todas las monedas, indicar qué monedas debe ir eligiendo Sophia para si misma y para Mateo, de tal forma que se asegure de ganar siempre. Considerar que Sophia siempre comienza (para sí misma).

\subsection{Planteamiento}

\subsection{Confirmar la optimalidad}